\section{引言与内容概览}

\subsection{讲座定位与学习目标}

Lecture 15 由 Chelsea Finn 深度解析层次化模仿与强化学习的工程框架和治理路径。本章目的不仅是让我们理解为什么层次结构在长视野任务中不可或缺，更是去复现讲者在讲座中不断强调的“结构化探查 → 观察 → 治理”闭环。讲座横跨数据、架构、prompt/目标设计和实践案例，整个笔记按教学逻辑重组，方便日后复现。

\begin{knowledgebox}{构建二层层次笔记的原则}
将讲座拆解为：1)动机、2)框架、3)子目标/模仿、4)语言+LLM调用、5)工程治理。每一块都需要用 `importantbox/knowledgebox/warningbox` 强调洞察，用表格或图片整理对比，最后用小结归纳以便利复盘。
\end{knowledgebox}

\subsection{时间线与素材地图}

\begin{table}[H]
\centering
\begin{tabular}{p{3.5cm}p{8cm}}
\toprule
时间区间 & 内容亮点与素材 \\
\midrule
00:00--00:12 & 讲者提出层次化原因、long horizon 与 hierarchical induction bias（slide-01～05） \\
00:12--00:30 & Options/SMDP 形式化与 value estimation（slide-06～12） \\
00:30--00:45 & Goal-conditioned 分层、HIRO/HAC 对比（slide-13～20） \\
00:45--00:55 & 层次化模仿与技能发现（slide-21～28） \\
00:55--01:12 & LLM+语言规划、高级案例与部署建议（slide-29～40） \\
01:12--01:20 & 实践 checklists + Q\&A 收官（slide-41～46） \\
\bottomrule
\end{tabular}
\caption{Lecture 15 时间线与 slide 对应。}
\end{table}

\begin{figure}[H]
\centering
\includegraphics[width=0.85\textwidth]{slides-images/slide-01.jpg}
\caption{讲座开场 slide：层次结构为何值得关注（slide-01）。}
\end{figure}

\begin{importantbox}{素材使用策略}
优先引用 slides 截图；若 slide 内容不足，再用字幕/笔记扩展。每一幅图像必须铺设目录中的路径（如 `slides-images/slide-xx.jpg`），并在 caption 里列出 timecode 便于查回。
\end{importantbox}

\subsection{本章小结}

本节通过定位、时间线和素材指引，构建了复盘本讲的地图。后续所有内容都依序围绕动机 → 模型 → 实践进行延展。

